\documentclass[aspectratio=169, 12pt]{beamer}

% Page geometry matching PPTX Widescreen (33.86cm x 19.05cm)
\geometry{paperwidth=33.86cm, paperheight=19.05cm}
\paperwidth=33.86cm
\paperheight=19.05cm
\pdfpagewidth=\paperwidth
\pdfpageheight=\paperheight

\usepackage{templates/icesibeamer/icesibeamer}
\usepackage{fontspec}
\usepackage[spanish]{babel}

% Institutional Font: Plus Jakarta Sans from resources/fonts/
\setmainfont{PlusJakartaSans-Regular.ttf}[
    Path = resources/fonts/,
    BoldFont = PlusJakartaSans-Bold.ttf
]
\setsansfont{PlusJakartaSans-Regular.ttf}[
    Path = resources/fonts/,
    BoldFont = PlusJakartaSans-Bold.ttf
]

\begin{document}

% SLIDE 1: Portada A
\titleSlideA{Título de la Diapositiva 1}{Subtítulo de la presentación}

% SLIDE 2: Portada B
\titleSlideB{Título de la Diapositiva 2}{Subtítulo de la presentación}{Autor / Nombre del evento / Fecha}

% SLIDE 3: Portada C
\titleSlideC{Título de la Diapositiva 3}{Subtítulo de la presentación}

% SLIDE 4: Portada D
\titleSlideD{Título de la Diapositiva 4}{Subtítulo de la presentación}{Texto Destacado en Caja Azul}

% SLIDE 5: Portada E
\titleSlideE{Título de la Diapositiva 5}{Subtítulo de la presentación}

% SLIDE 6: Portada F
\titleSlideF{Título de la Diapositiva 6}{Subtítulo de la presentación}

% SLIDE 7: División A
\sectionSlideA{División A: Título de Sección}

% SLIDE 8: División B
\sectionSlideB{División B: Título de Sección}{resources/photo_section_b.jpg}

% SLIDE 9: División C
\sectionSlideC{División C: Título de Sección}

% SLIDE 10: División E (5 variantes de color)
\sectionSlideEBlue{División E: Fondo Azul}
\sectionSlideE{División E: Fondo Blanco}
\sectionSlideEGreen{División E: Fondo Verde}
\sectionSlideEYellow{División E: Fondo Amarillo}
\sectionSlideEOrange{División E: Fondo Naranja}

% SLIDE 15: Sidebar Izquierda Naranja
\slideSidebarLeftOrange{Título Diapositiva 15}{Esta es una diapositiva con una barra lateral izquierda de color naranja y el logotipo institucional en blanco negativo. El contenido y el título se posicionan limpiamente en la parte derecha.}

% SLIDE 16: Sidebar Izquierda Azul
\slideSidebarLeftBlue{Título Diapositiva 16}{Esta variante cuenta con la barra lateral en color azul institucional, manteniendo el mismo esquema de alineación y márgenes de seguridad.}

% SLIDE 17: Sidebar Izquierda Morada
\slideSidebarLeftPurple{Título Diapositiva 17}{Esta variante cuenta con la barra lateral en color morado, ideal para cambiar de sección de forma dinámica.}

% SLIDE 18: Franja Superior con Acento a la Izquierda
\slideStripeTopLeft{Título Diapositiva 18}{Esta es una diapositiva con una franja azul superior y una franja verde de acento posicionada a la izquierda. La caja de contenido se ubica en el centro de la diapositiva sobre fondo blanco.}

% SLIDE 19: Franja Superior con Acento a la Derecha
\slideStripeTopRight{Título Diapositiva 19}{Esta es una diapositiva hermana con una franja azul superior y la franja verde de acento posicionada a la derecha. Ambas comparten la alineación y el uso del logo blanco en negativo en la esquina superior izquierda.}
\end{document}